\chapter{Core Vocabulary}
\label{ch:vocabulary}

\epigraph{Terminology is prevented from becoming detached from operational meaning.
The insistence on repeatedly rebuilding assumptions from the ground upward functions
as a defense against proxy drift at the conceptual level itself.}{}

\noindent This chapter establishes the definitional scaffold on which the remainder of the
monograph is built. Terms introduced here are used without redefinition in subsequent
chapters. Where a concept is partially anticipated in Chapter~\ref{ch:manifold}, the
definition here provides the formal completion.

\section{State Spaces and Trajectory Selection}

\begin{definition}{State Space}
A \emph{state space} $\mathcal{X}$ is a set whose elements represent all possible
configurations of a system. A \emph{trajectory} through $\mathcal{X}$ is a map
$\gamma: I \to \mathcal{X}$ from an index set $I$ (typically a time interval) into
the state space. The \emph{selection pressure} acting on a trajectory space is the
set of mechanisms that filter admissible from inadmissible paths. Selection pressure
may arise from physical laws, thermodynamic constraints, logical requirements,
historical boundary conditions, or any combination of these.
\end{definition}

The distinction between exhaustive search and constrained navigation is fundamental.
A system performing exhaustive search considers the full space $\mathcal{X}^I$ of
possible trajectories and evaluates each against an objective function. A system
performing constrained navigation operates only within $\mathcal{A}(\mathcal{C})$
from the outset, never considering inadmissible paths. Biological and cognitive systems
almost exclusively employ constrained navigation, because exhaustive search is
thermodynamically impossible at any scale of practical interest.

\section{Compression, Sparsity, and Structural Invariants}

\begin{definition}{Compression}
A \emph{compression} of a system description $D$ is a map $\kappa: D \to D'$ where
$|D'| < |D|$ under some measure of description length, and $D'$ preserves a specified
set of structural relations present in $D$. A compression is \emph{lossless with
respect to} $\mathcal{R}$ if every relation $r \in \mathcal{R}$ is recoverable from
$D'$. It is \emph{lossy} if some $r \in \mathcal{R}$ is not recoverable. Compression
is therefore not reduction of volume alone but selective preservation of relations.
\end{definition}

\begin{definition}{Structural Invariant}
A \emph{structural invariant} of a system under a family of transformations
$\{T_\alpha\}$ is a property $P$ such that $P(T_\alpha(s)) = P(s)$ for all $\alpha$
and all states $s$ in the relevant domain. Structural invariants are the properties
that compression must preserve to remain operationally useful. A compression that
erases structural invariants necessary for understanding, provenance, or accountability
is a \emph{destructive compression} regardless of its information-theoretic efficiency.
\end{definition}

\begin{definition}{Sparsity}
A representation is \emph{sparse} if most of its degrees of freedom are zero or
near-zero, with significant values concentrated in a small fraction of components.
Sparsity is the condition under which compression becomes tractable: a sparse
representation can be compressed without destroying structural invariants because
most of the space is already empty. The \emph{Natural Sparsity Principle} states
that in biological systems, sparsity emerges as a natural consequence of
physiological pressures---energetic constraints, signal noise, chemical
gradients---without requiring explicit computational penalties to enforce minimal
complexity.
\end{definition}

\subsection*{Why Sparsity Makes Reconstruction Tractable}

The connection between sparsity and tractable reconstruction can be sketched
formally. Let $\mathbf{x} \in \mathbb{R}^N$ be a signal with $k \ll N$ nonzero
components. The number of bits required to specify $\mathbf{x}$ is approximately
$k \log(N/k) + k \log(\max|\mathbf{x}|/\epsilon)$ for precision $\epsilon$, far
less than the $N \log(\max|\mathbf{x}|/\epsilon)$ bits required for a dense signal.

Under bounded energetic constraints, a system with budget $B$ (in bits or equivalent
energy units) can faithfully represent $\mathbf{x}$ if and only if $B \geq$ the
encoding cost. For dense signals this budget is exhausted rapidly; for sparse signals
it is not. The compressive sensing literature shows that under sparsity, signals can
be recovered from $O(k \log N)$ measurements rather than $N$ measurements, with
reconstruction error bounded by the sparsity level.

The biological relevance is direct. Neural firing rates are sparse: typical cortical
neurons fire at low average rates, with high-rate bursts carrying most information.
This means the energetic cost of neural representation scales with signal sparsity
rather than state-space dimensionality. A brain that maintained dense representations
of all sensory dimensions simultaneously would require thermodynamically infeasible
metabolic support. Sparsity is not a design choice but a thermodynamic necessity
that the Natural Sparsity Principle identifies as a universal constraint on
cognitive architecture.

\section{Reconstruction and Hallucination}

\begin{definition}{Reconstruction}
\emph{Reconstruction} is the process by which a system assembles an operationally
coherent configuration from partial traces. A reconstruction is \emph{grounded} if
it remains constrained by external feedback signals. It is \emph{corrigible} if
errors in reconstruction are detectable and correctable through interaction with the
environment. It is \emph{drifting} if it progressively decouples from external
constraint and stabilizes on an internally self-consistent but externally inaccurate
model.
\end{definition}

\begin{definition}{Hallucination (Formal)}
\emph{Hallucination} is reconstruction that exceeds available constraint: the
assembly of a locally coherent configuration from traces that underdetermine it,
without external correction signals sufficient to resolve the underdetermination.
Hallucination is not a pathology but the default operating condition of any
sufficiently compressed system under uncertainty. The critical distinction is not
between hallucinating and non-hallucinating systems but between systems whose
hallucinations remain \emph{corrigible} through feedback and systems whose internal
reconstructions drift irreversibly from external constraint.
\end{definition}

\begin{remark}
This definition of hallucination is deliberately broader than its use in AI
discourse, where it typically refers to language model confabulation. The formal
definition encompasses perception, memory consolidation, theoretical inference,
institutional modeling, and any other process in which a partial trace is assembled
into a more complete representation. All such processes hallucinate in this sense;
the question is always whether the hallucination is correctable.
\end{remark}

\subsection*{Reconstruction Bound}

The inevitability of hallucination under compression can be made precise. Let $D$
be a description of cardinality $|D| = N$ and let $\kappa(D) = D'$ be a compression
with $|D'| = M < N$. The number of distinct descriptions $D$ consistent with $D'$
is at least $2^{N-M}$, since the compressed representation cannot distinguish among
them. Any reconstruction process that must select a single completion from among
these $2^{N-M}$ possibilities must hallucinate: it must choose structure that the
compression does not uniquely determine.

More formally, if $p(D \mid D')$ is the posterior distribution over original
descriptions given the compressed form, then:
\[
  H(D \mid D') = -\sum_D p(D \mid D') \log p(D \mid D') \geq N - M > 0.
\]
The conditional entropy is strictly positive whenever $M < N$. A reconstructing
system with access only to $D'$ cannot recover $D$ with certainty. It must
interpolate the missing $N - M$ bits from prior structure, contextual signals,
or statistical regularities. That interpolation is hallucination in the formal
sense: reconstruction exceeding available constraint.

The corrigibility condition---the requirement that reconstruction errors remain
correctable through feedback---does not eliminate this hallucination. It ensures
only that when the interpolated structure turns out to be wrong, the error is
detectable and correctable before it propagates into downstream commitments.

\section{Provenance and the Boot Sequence}

\begin{definition}{Provenance}
The \emph{provenance} of a conceptual or computational object is the recoverable
dependency graph of the operations and inputs that produced it. Provenance includes:
the sequence of constraint-stabilization events that enabled the object's current
form; the assumptions that were imported versus derived; and the points at which
external evidence entered the reconstruction process. A system with high provenance
transparency is one whose dependency graph can be approximately reconstructed by an
external observer.
\end{definition}

\begin{definition}{Boot Sequence}
A \emph{boot sequence} is the ordered set of constraint-stabilization events
necessary to make a higher-level abstraction operationally meaningful. Lower-level
constraints must be stabilized before higher-level ones can be evaluated. Erasing
the boot sequence---presenting only the terminal abstraction without the ordered
initialization history---transforms an inspectable object into an opaque one. The
boot sequence is thus the provenance of a system's layered structure.
\end{definition}

\section{The Formal Substrates: RSVP, KES, Spherepop}

Three formal frameworks function as computational and geometric substrates throughout
this monograph. They are introduced briefly here and developed formally in the
appendices. References to these frameworks in subsequent chapters should be understood
as references to the formal specifications given there.

\begin{definition}{\RSVP Field System}
The \emph{Relativistic Scalar-Vector Plenum} (\RSVP) is a coupled field system
consisting of three dynamical fields over a spacetime manifold $\mathcal{M}$:
\begin{itemize}
  \item $\Phifield(\mathbf{x}, t)$: a scalar density field encoding concentration,
        potential, or legitimacy of state.
  \item $\vfield(\mathbf{x}, t)$: a vector flow field encoding directed transport,
        trajectories, and causality.
  \item $\Sfield(\mathbf{x}, t)$: an entropy field encoding disorder, uncertainty,
        and complexity budgets.
\end{itemize}
The three fields evolve according to coupled nonlinear partial differential equations.
The full Lagrangian structure and seven governing axioms are given in
Appendix~\ref{app:rsvp}. Within the semantic interpretation developed in
Chapter~\ref{ch:yarncrawler_formal}, $\Phifield$ measures semantic density, $\vfield$
encodes recursive trajectories of meaning, and $\Sfield$ tracks unresolved noise
available for future reinterpretation.
\end{definition}

\begin{definition}{\KES Map}
The \emph{Kinetic-Event Synthesis} (\KES) map is the function
\[
  \KES: \Omega_t \longrightarrow H_{t+1}
\]
from a current world-state $\Omega_t$ to a successor hypothesis space $H_{t+1}$.
The map formalizes the basic cognitive operation: given a current configuration of
constraints, evidence, and prior commitments, what hypothesis spaces are admissible
at the next step? The full convergence conditions and technical specification are
given in Appendix~\ref{app:kes}.
\end{definition}

\begin{definition}{Spherepop Calculus}
\emph{Spherepop} is an irreversible event calculus defined by four primitive
operators acting on a possibility space $\mathcal{P}$:
\begin{itemize}
  \item $\mathrm{Pop}$: extracts committed events from possibility space,
        transitioning them to actuality.
  \item $\mathrm{Refuse}$: eliminates inadmissible trajectories, contracting the
        admissible space.
  \item $\mathrm{Bind}$: establishes structural dependencies between events,
        constraining future possibilities.
  \item $\mathrm{Collapse}$: commits a possibility space to a definite
        configuration, rendering alternatives inaccessible.
\end{itemize}
All four operations are irreversible: once applied, they cannot be undone within
the same system. The operational semantics of all four operators are developed
formally in Appendix~\ref{app:spherepop}.
\end{definition}

\begin{remark}
The irreversibility of Spherepop operators is not a limitation but a design
commitment. Systems that can always reverse their operations accumulate no
structural residue; they have no history and therefore no trajectory. The
Collapse operator in particular formalizes the commitment structure that
distinguishes genuine decision from perpetual deferral.
\end{remark}

\section*{Constraint Boundary}

\begin{constraintboundary}
\textbf{What this chapter establishes:} Formal definitions of admissibility,
compression, sparsity, reconstruction, hallucination, provenance, and the boot
sequence. The reconstruction bound $H(D \mid D') \geq N - M > 0$ is formally
derived. The three formal substrates (\RSVP, \KES, Spherepop) are introduced as
technical objects.

\textbf{What this chapter does not establish:} That these definitions uniquely
capture the relevant phenomena in any particular domain. The definitions are
productive insofar as they generate useful inferences; their domain of validity
must be established by application.

\textbf{Remaining speculative:} The claim that the Natural Sparsity Principle
operates universally across biological systems is a \textbf{[C]} conjecture
supported by neuroscience and thermodynamics but not formally derived here.

\textbf{Falsification conditions:} A system that produces corrigible reconstruction
from a strictly lossless compression would falsify the claim that hallucination is
constitutive of compressed inference. A biological cognitive system that maintains
dense rather than sparse representations at thermodynamically sustainable cost
would falsify the Natural Sparsity Principle.
\end{constraintboundary}
